\documentclass[12pt,a4paper]{article}
\usepackage[T2A]{fontenc}
\usepackage[utf8]{inputenc}
\usepackage[russian,english]{babel}
\usepackage{amsmath,amssymb}
\usepackage[margin=2.5cm]{geometry}
\usepackage[expansion=false]{microtype}
\DisableLigatures{encoding=*,family=*}
\setlength{\parindent}{0pt}
\renewcommand{\baselinestretch}{1.4}

\begin{document}
\selectlanguage{russian}
\thispagestyle{empty}

% ---------- Журнальная шапка ----------
\begin{flushleft}\small
\textit{В редакцию журнала \guillemotleft Программирование\guillemotright{} (Programming and Computer Software), ИСП РАН}\\[2pt]
Рубрика: Языки, компиляторы, системы программирования и аппаратно-зависимые вычисления\\[2pt]
УДК~004.222
\end{flushleft}

\vspace{0.4cm}
\hrule
\vspace{0.6cm}

% ---------- Русский блок ----------
\begin{center}
{\large\textbf{GoldenFloat: семейство форматов с плавающей точкой\\
со статическим разбиением, выведенным из золотого сечения~\(\varphi\),\\
от GF4 до GF256, с точным целочисленным тождеством Люка}}\\[10pt]
\textbf{\copyright~2026~г.\quad Д.\,В.~Васильев}\\[6pt]
\textit{Независимый исследователь (Trinity~S\textsuperscript{3}AI), Ко~Самуи, Королевство Таиланд}\\[4pt]
e-mail: admin@t27.ai \quad\textbar\quad ORCID: 0009-0008-4294-6159
\end{center}

\vspace{0.3cm}
\begin{flushleft}\small
Поступила в редакцию: \rule{2.8cm}{0.4pt}\quad
После доработки: \rule{2.8cm}{0.4pt}\quad
Принята к публикации: \rule{2.8cm}{0.4pt}\\[4pt]
\textbf{DOI:}~\rule{4.5cm}{0.4pt}\quad\textit{(присваивается редакцией и издательством)}
\end{flushleft}

\vspace{0.3cm}
\textbf{Аннотация.}
Дано аппаратно-ориентированное описание семейства числовых форматов GoldenFloat~(GF), порождаемого единым замкнутым правилом разбиения разрядной сетки: для каждой полной разрядности \(N\ge 4\) ширина поля порядка равна \(e=\mathrm{round}((N-1)/\varphi^{2})\), ширина мантиссы \(f=N-1-e\), где \(\varphi=(1+\sqrt5)/2\). Представлены три воспроизводимых артефакта: (i)~открытый многоразрядный RTL-генератор (GF4--GF256) с дифференциальной прогонкой в режиме непрерывной интеграции относительно корректно округлённого эталона; (ii)~целочисленный тракт аккумулятора, точный по числам Люка, проверенный с точностью до 500~знаков для \(n=1,\dots,256\); (iii)~кодек GF16 на FPGA, проходящий тестовое окружение 35/35 на частоте 323~МГц (Artix-7, Xilinx~XC7A35T). Правило воспроизводит реализованные ширины порядка девяти форматов (9/9) и согласованно продолжается на бо\'{}льшие разрядности. Работа позиционируется рядом с posit (Posit~Standard~2022), takum, OCP-MX и черновиком IEEE~P3109; тезис о широте/согласованности зафиксирован как открытая гипотеза с предрегистрированным путём фальсификации (реестр~FL-002), без заявлений о превосходстве по точности.\\[6pt]
\textbf{Ключевые слова:} численные форматы, форматы с плавающей точкой, золотое сечение, числа Люка, posit, takum, OCP-MX, аппаратная арифметика.

\vspace{0.5cm}
\hrule
\vspace{0.6cm}

% ---------- Английский блок ----------
\selectlanguage{english}
\begin{center}
{\large\textbf{GoldenFloat: a family of static-split floating-point formats\\
derived from the golden ratio \(\varphi\), from GF4 to GF256,\\
with an exact integer Lucas identity}}\\[10pt]
\textbf{\copyright~2026\quad D.\,V.~Vasilev}\\[6pt]
\textit{Independent researcher (Trinity~S\textsuperscript{3}AI), Ko~Samui, Kingdom of Thailand}\\[4pt]
e-mail: admin@t27.ai \quad\textbar\quad ORCID: 0009-0008-4294-6159
\end{center}

\vspace{0.3cm}
\textbf{Abstract.}
We present a hardware-oriented description of GoldenFloat~(GF), a family of static-split floating-point formats generated by a single closed-form rule: for every full width \(N\ge 4\) the exponent-field width is \(e=\mathrm{round}((N-1)/\varphi^{2})\) with fraction width \(f=N-1-e\) and \(\varphi=(1+\sqrt5)/2\). Three reproducible artefacts are given: (i)~an open multi-width RTL generator (GF4--GF256) with a continuous-integration differential sweep against a correctly-rounded reference; (ii)~a Lucas-exact integer accumulator path verified to 500 digits for \(n=1,\dots,256\); (iii)~a GF16 FPGA codec passing a 35/35 testbench at 323~MHz (Artix-7, Xilinx~XC7A35T). The rule reproduces the realised exponent widths of nine formats (9/9) and extends consistently to further widths. The work is positioned alongside posit (Posit Standard 2022), takum, OCP-MX and the IEEE~P3109 draft; the breadth/coherence thesis is registered as an open conjecture with a pre-registered falsification path (ledger~FL-002), with no superiority claim.\\[6pt]
\textbf{Keywords:} numeric formats; floating point; golden ratio; Lucas numbers; posit; takum; OCP-MX; arithmetic hardware.

\vspace{0.6cm}
\hrule
\vspace{0.4cm}
\selectlanguage{russian}
\small
\textbf{Сведения для редакции.} Статья оригинальна, ранее не публиковалась и не находится на рассмотрении в других изданиях. Конфликт интересов отсутствует, внешнее финансирование не привлекалось. Автор согласен на двойное слепое рецензирование. Как гражданин~РФ, по запросу редакции готов предоставить заключение о возможности открытого опубликования. Объём статьи 24~стр.; подготовлена в \LaTeX{} (прилагается исходный файл \texttt{.tex} и собранный PDF).

\end{document}
